\subsection{Geometric oracle and identifiability}
\label{sec:oracle}

\textbf{Structures and labels.} A crystal structure is a triple $s = (L, X, Z)$: a lattice
$L \in \mathbb{R}^{3\times3}$, fractional coordinates $X = \{x_i\}_{i=1}^{N} \subset [0,1)^3$, and species
$Z = \{z_i\}$. The label map is $y(s) = \spg(s)$, the crystal system returned by a deterministic symmetry
algorithm (spglib) at a fixed tolerance $\tau$, taking values in the seven crystal systems. Because the
label is computed rather than annotated, label correctness is not a competing explanation for model error;
Section~\ref{sec:benchmark} describes the tolerance sweep that certifies it.

\textbf{Renders.} The frozen protocol renders a $2\times2\times2$ tiling of the conventional cell through
five orthographic cameras $\mathcal{C} = \{c_v\}_{v=1}^{5}$ (three principal-axis and two oblique views) at
$768$ px, giving images $I_v = \mathcal{R}(s, c_v)$. The tiling is a legibility choice inherited from the
render convention and is not information-free; Section~\ref{sec:conventions} measures what it costs. Each camera is a known unit direction $d_v$ with orthographic projection
$\Pi_v : \mathbb{R}^3 \to \mathbb{R}^2$, $\Pi_v(p) = P_v\, p$ for a fixed $2\times3$ matrix $P_v$ with
orthonormal rows perpendicular to $d_v$. The task given to a model is: from $(I_1,\dots,I_5)$, name the
crystal system.

\textbf{The oracle.} The oracle answers the same question from the renders' own geometry instead of from
pixels. It forward-projects the ground-truth atom positions of the conventional cell through each frozen
camera, \emph{discards} the correspondence between views, re-solves it by ray intersection with cross-view
verification, and applies $\spg$ to the reconstruction (Algorithm~\ref{alg:oracle}). It is deterministic
and CPU-only. Two properties of that definition are load-bearing and we state them here rather than in the
limitations. First, the oracle assumes \emph{ideal extraction}: every atom's centroid is available in every
view, including atoms an image hides behind another, so $R_1$ bounds identifiability from the projection
geometry and not from the pixels. A visibility-conditioned variant that removes occluded detections is a
separate measurement and is reported in Section~\ref{sec:conventions}. Second, the oracle inverts the
conventional-cell projections rather than the tiled ones the renderer draws; supplying it the tiled geometry
instead is one of the interventions of Section~\ref{sec:conventions} and costs it twenty-two structures.

\begin{algorithm}[t]
\caption{Geometric oracle $O(s, \mathcal{C})$}
\label{alg:oracle}
\begin{algorithmic}[1]
\For{each view $v$}
  \State $Q_v \gets \{(\Pi_v(L x_i),\, z_i) : i = 1,\dots,N\}$
  \Comment{projected, species-labelled 2D points; cross-view identity discarded}
\EndFor
\For{each pair of views $(v, w)$ and each same-species pair $(q, q') \in Q_v \times Q_w$}
  \State $\hat{p} \gets$ intersection of the back-projected rays $\Pi_v^{-1}(q)$ and $\Pi_w^{-1}(q')$
  \State accept $\hat{p}$ iff its projection lies within tolerance of a same-species point in every
         remaining view
\EndFor
\State \Return $\spg(\text{verified reconstruction})$
\end{algorithmic}
\end{algorithm}

\begin{definition}[Identifiability ceiling]
\label{def:ceiling}
For an evaluation sample $\mathcal{S}$, the ceiling is
\[
R_1 = \frac{1}{|\mathcal{S}|}\sum_{s \in \mathcal{S}} \mathbf{1}\!\left[\,O(s,\mathcal{C}) = y(s)\,\right],
\]
the fraction of structures whose label the frozen renders determine for an ideal geometric reader.
\end{definition}

\begin{proposition}[Exact recovery]
\label{prop:ident}
Let $s$ be a structure rendered under known orthographic cameras. Suppose (i) at least two of the camera
directions are linearly independent, and (ii) the cross-view verification of
Algorithm~\ref{alg:oracle} admits a unique consistent correspondence, i.e.\ no two distinct atoms produce
verified ray intersections within the matching tolerance of one another. Then the oracle recovers every
atom position exactly, and $O(s, \mathcal{C}) = \spg(s) = y(s)$.
\end{proposition}

A proof sketch is in Appendix~\ref{app:proof}; the argument is elementary (two orthographic images along
independent directions determine a point as the intersection of two affine lines, and verification against
the remaining views eliminates spurious matches outside the coincidence set). The proposition is useful for
what its failure set says. Condition (i) holds by construction for the frozen camera set. Condition
(ii) fails under projective coincidence, when distinct atoms project onto near-identical pixels, and this
is exactly where the measured ceiling falls short of $1$: coincidence grows with atoms per cell, which
predicts the cell-size dependence observed at scale-up (Section~\ref{sec:controls}). A separate quantity
governs what an extractor reading pixels could recover: across the frozen five-view set, the fraction of
atoms with fewer than two clear views is 0.26\% (original sample) and 0.87\% (replication sample), so the
protocol's occlusion leaves almost every atom triangulable in principle (Section~\ref{sec:conventions}).

\textbf{Why a ceiling below $1.0$ is the point.} A design that supplied the source coordinates as text
would have a ceiling of exactly $1.0$ by construction and would measure nothing about the images. A ceiling
below $1.0$ says the protocol withholds information, and Proposition~\ref{prop:ident} locates where.
